\scanpage{81}
% Faint hooked trace at the upper left edge of the sheet.
\[
H^*(S,\operatorname{End}(M))\mathrel{\unsure{\relbar|\longrightarrow}}
H^*(S,\operatorname{End}(M)\otimes_{\Lambda}\underline{O}_S)\simeq H^{\unsure{2|*}}(S,\operatorname{End}(E))
\]

Enfin, sous \unsure{les conditions|la condition} de \unsure{1.5|(1.5}, $M_0$ \unsure{etant|étant} aussi plat sur $\Lambda$, on aura que $\unsure{\xi|\zeta}_{\Lambda}$ est l'image d'un élément
\[
\xi_0\in H^2(B_{\pi},\operatorname{End}(M_0))\quad .
\]

(NB Pour calculer $\xi_0$, par fonctorialité en $\pi$, on est ramené aussitôt au cas ``universel'' $\pi=\operatorname{Aut}(M_0)$, et on peut \unsure{évidemment|évidement} supposer $M_0=\Lambda^d=(\mathbb Z/n\mathbb Z)^d$, donc $\pi=\operatorname{Gl}(d,\mathbb Z/n\mathbb Z)$....) D'après \uline{Illusie}, c'est l'obstruction à remonter la représentation de $\pi$ dans $M_0$ en une représentation dans un module libre sur $\mathbb Z/\unsure{n|p}^2\mathbb Z$. \inserted{Lorsque $n=p$ (premier)} Elle est \unsure{évidemment|évidement} nulle si $M_0$ est de rang $\leq 1$, mais sans doute pas nulle si $M_0$ de rang $d\geq 2$, dans le cas universel $\pi=\deleted{\unsure{A}}\operatorname{Gl}(d,\mathbb F_p)$.

Notons que si le rev. $S'$ de $S$ est le rev. universel, alors
\[
H^2(B_{\pi},-)\longrightarrow H^2(S,-)
\]
est \uline{injectif}, donc on \unsure{trouvera|trouve} un élément non nul \deleted{dans} $\xi_{\Lambda}\in H^2(S,\operatorname{End}(M))$.
% archive sheet number: 80
